%% main_report61.tex
%% SAMBP Technical Report TR-61/2026
%% Relay 64G Stator Earth Fault Protection: Optimal Sub-harmonic Injection
%% Frequency under IBR Harmonic Contamination
%% =========================================================================

\documentclass[11pt,a4paper]{article}

\usepackage[T1]{fontenc}
\usepackage[utf8]{inputenc}
\usepackage{lmodern}
\usepackage[margin=25mm]{geometry}
\usepackage{amsmath,amssymb,amsthm}
\usepackage{booktabs}
\usepackage{array}
\usepackage{longtable}
\usepackage{graphicx}
\usepackage{xcolor}
\usepackage[hidelinks]{hyperref}
\usepackage{microtype}
\usepackage{pgfplots}
\pgfplotsset{compat=1.18}
\usepackage{tikz}
\usetikzlibrary{arrows.meta,positioning}
\usepackage{siunitx}
\usepackage{cite}

\definecolor{sambpblue}{RGB}{0,70,140}
\definecolor{darkgreen}{RGB}{0,120,50}
\definecolor{warningred}{RGB}{180,0,0}

\newtheorem{proposition}{Proposition}
\newtheorem{remark}{Remark}

\newcommand{\Ihealthy}{I_{\mathrm{healthy}}}
\newcommand{\Ifault}{I_{\mathrm{fault}}}
\newcommand{\Itrip}{I_{\mathrm{trip}}}
\newcommand{\sigman}{\sigma_n}
\newcommand{\fopt}{f^*}
\newcommand{\fsub}{f_{\mathrm{sub}}}
\newcommand{\fconv}{f_{\mathrm{conv}}}
\newcommand{\RN}{R_N}
\newcommand{\CG}{C_G}

\title{%
  \textbf{\textsf{\color{sambpblue}SAMBP Technical Report TR-61/2026}}\\[6pt]
  \Large Relay 64G Stator Earth Fault Protection:\\
  Optimal Sub-harmonic Injection Frequency under\\
  IBR Harmonic Contamination%
}
\author{SAMBP Research Group}
\date{April 2026 \quad\textbar\quad Chapter 6, \S6.8 \quad\textbar\quad
      Target: \emph{IEEE Transactions on Industrial Electronics}}

\begin{document}
\maketitle
\thispagestyle{empty}
\tableofcontents
\clearpage

%% ────────────────────────────────────────────────────────────────────────────
\section{Background and Problem Statement}
\label{sec:background}
%% ────────────────────────────────────────────────────────────────────────────

Stator earth fault protection (Relay 64G) in generators relies on sub-harmonic
voltage injection at the machine neutral.  A voltage source $V_{\rm inj}$ is
applied between the neutral and ground via a neutral grounding resistor $\RN$,
and the return current is monitored.  A fault at any point on the stator
winding creates a low-impedance conductive path to ground, increasing the
measured current above a threshold $\Itrip$.

Standard practice uses a fixed injection frequency $\fconv = 20$\,Hz, well
below the rated 50\,Hz to avoid interference with the fundamental.  However,
high-penetration inverter-based resources (IBR) generate sub-harmonic
currents in the range 15--25\,Hz due to:
\begin{itemize}
  \item Converter dead-time compensation harmonics,
  \item LCL-filter resonance under weak-grid conditions,
  \item GFM low-frequency droop dynamics.
\end{itemize}
When the IBR noise peak $\fsub \approx \fconv = 20$\,Hz, the standard relay
suffers simultaneous \emph{security} and \emph{dependability} failures:

\begin{description}
  \item[Security failure (false alarm).]
    IBR noise adds to the measured healthy current $\Ihealthy(\fconv)$.
    If $\Ihealthy(\fconv) + 2\sigman(\fconv) > \Itrip$, a false trip occurs.
  \item[Dependability failure (missed detection).]
    IBR noise subtracts from the fault current measurement.  For high-resistance
    faults with $\Ifault \approx \Itrip$, noise corrupts the safety margin and
    causes missed detections.
\end{description}

This report proposes selecting the injection frequency adaptively:
\begin{equation}
  \fopt = \arg\max_{f \in \mathcal{F}} \min\bigl(
    \underbrace{\Itrip - \Ihealthy(f) - 2\sigman(f)}_{\text{security margin}},\;
    \underbrace{\Ifault(R_f) - \Itrip - 2\sigman(f)}_{\text{dependability margin}}
  \bigr)
  \label{eq:opt}
\end{equation}
where $\mathcal{F} = [20, 40]$\,Hz is the permissible injection range (below the
fundamental, above converter switching noise).

%% ────────────────────────────────────────────────────────────────────────────
\section{Physical Model}
\label{sec:model}
%% ────────────────────────────────────────────────────────────────────────────

\subsection{Healthy winding current}
\label{sec:healthy}

Without a fault, the injection source drives a displacement current through the
distributed stator-to-ground capacitance $\CG$:
\begin{equation}
  \Ihealthy(f) = \frac{V_{\rm inj} \cdot \omega \CG}{\sqrt{1 + (\omega \CG \RN)^2}},
  \qquad \omega = 2\pi f.
  \label{eq:Ihealthy}
\end{equation}
This current \emph{increases monotonically with} $f$ (as capacitive reactance
$X_c = 1/(\omega\CG)$ decreases), eventually exceeding $\Itrip$ and setting a
hard upper bound on permissible $\fopt$.

\subsection{Fault current}
\label{sec:fault}

A stator earth fault at winding position $\theta$ (fraction from neutral)
creates a direct conductive path.  At sub-harmonic frequencies the winding
impedance from neutral to $\theta$ is negligible, so:
\begin{equation}
  \Ifault(R_f) = \frac{V_{\rm inj}}{\RN + R_f},
  \label{eq:Ifault}
\end{equation}
which is \emph{frequency-independent}.  The fault resistance $R_f$ is the
dominant variable; high-resistance faults ($R_f \to 1/\Itrip - \RN$) define the
protection sensitivity limit.

\begin{remark}[Position independence]
  The conductive fault current is independent of winding position $\theta$,
  providing 100\% winding coverage at all sub-harmonic injection frequencies.
  This is a key advantage of neutral injection over zero-sequence voltage
  methods.
\end{remark}

\subsection{IBR noise model}
\label{sec:noise}

The IBR sub-harmonic noise at the neutral terminal is modelled as a Gaussian
spectral peak:
\begin{equation}
  \sigman(f) = A_{\rm noise} \cdot \exp\!\left(-\frac{(f - \fsub)^2}{2\sigma_{\rm sub}^2}\right),
  \label{eq:noise}
\end{equation}
with peak frequency $\fsub \in [15, 25]$\,Hz, amplitude $A_{\rm noise}$, and
bandwidth $\sigma_{\rm sub} = 3$\,Hz.

\subsection{Detection criterion}
\label{sec:criterion}

Using a $2\sigma$ statistical margin:
\begin{align}
  \text{Security:}\quad &  \Ihealthy(f) + 2\sigman(f) < \Itrip
  \label{eq:security} \\
  \text{Dependability:}\quad & \Ifault(R_f) - 2\sigman(f) > \Itrip
  \label{eq:depend}
\end{align}

%% ────────────────────────────────────────────────────────────────────────────
\section{Optimal Frequency Proposition}
\label{sec:proposition}
%% ────────────────────────────────────────────────────────────────────────────

\begin{proposition}[Optimal frequency bounds]
  \label{prop:fopt}
  Let $\fsub \in [15, 25]$\,Hz and $A_{\rm noise} > 0$.  The optimal injection
  frequency $\fopt$ defined by \eqref{eq:opt} satisfies:
  \begin{equation}
    f_{\rm low} < \fopt \leq f_{\rm up}
    \label{eq:bounds}
  \end{equation}
  where $f_{\rm low} = \fsub + c_1 \sigma_{\rm sub}$ (above the noise peak by
  $c_1 \approx 2$--$3$ bandwidths) and $f_{\rm up}$ is defined by
  $\Ihealthy(f_{\rm up}) = \Itrip$ (security crossover frequency).  For the
  parameters in Table~\ref{tab:params}, $f_{\rm up} \approx 35$\,Hz, giving
  $\fopt \in [25, 35]$\,Hz for $\fsub \in [15, 25]$\,Hz.
\end{proposition}

\begin{proof}[Proof sketch]
  For $f < \fsub$: $\sigman(f) > \sigman(\fsub + (f-\fsub)) = \sigman(f')$
  for the mirror point $f' = 2\fsub - f > \fsub$.  So any $f < \fsub$ has
  a strictly worse noise level than the corresponding point above $\fsub$,
  with equal or worse healthy current (since $\Ihealthy$ is monotone).
  Hence $\fopt > \fsub$.  The upper bound follows directly from
  $\Ihealthy(f_{\rm up}) = \Itrip$ and the monotonicity of $\Ihealthy$.
\end{proof}

%% ────────────────────────────────────────────────────────────────────────────
\section{Validation Methodology}
\label{sec:validation}
%% ────────────────────────────────────────────────────────────────────────────

\subsection{System parameters}
\label{sec:params}

\begin{table}[h]
\centering
\caption{TR-61 system parameters}
\label{tab:params}
\begin{tabular}{llll}
  \toprule
  Parameter & Symbol & Value & Units \\
  \midrule
  Stator-to-ground capacitance & $\CG$   & $4.55\times10^{-4}$ & pu$\cdot$s \\
  Neutral grounding resistor   & $\RN$   & 0.080 & pu \\
  Injection voltage            & $V_{\rm inj}$ & 1.000 & pu \\
  Trip threshold impedance     & $Z_{\rm trip}$ & 10.0 & pu \\
  Trip current threshold       & $\Itrip$ & 0.100 & pu \\
  Noise bandwidth              & $\sigma_{\rm sub}$ & 3.0 & Hz \\
  Security crossover freq.     & $f_{\rm up}$  & $\approx$35 & Hz \\
  \bottomrule
\end{tabular}
\end{table}

With these parameters: $\Ihealthy(20\,\mathrm{Hz}) = 0.057$\,pu $\ll \Itrip$;
$\Ihealthy(35\,\mathrm{Hz}) \approx \Itrip$.  The minimum detectable fault
resistance is $R_{f,\rm max} = V_{\rm inj}/\Itrip - \RN = 9.92$\,pu.

\subsection{Deterministic scenario matrix}

\begin{center}
\begin{tabular}{lllll}
  \toprule
  ID & Type & $R_f$ (pu) & $\fsub$ (Hz) & $A_{\rm noise}$ (pu) \\
  \midrule
  S01 & Security & — & 20.0 & 0.080 \\
  S02 & Security & — & 22.0 & 0.080 \\
  S03 & Security & — & 18.0 & 0.080 \\
  S04 & Security & — & 23.0 & 0.070 \\
  S05 & Security & — & 20.0 & 0.100 \\
  S06 & Fault & 6.0 & 20.0 & 0.080 \\
  S07 & Fault & 7.0 & 20.0 & 0.080 \\
  S08 & Fault & 7.0 & 22.0 & 0.080 \\
  S09 & Fault & 6.5 & 18.0 & 0.060 \\
  S10 & Fault & 7.0 & 20.0 & 0.090 \\
  \bottomrule
\end{tabular}
\end{center}

Security scenarios (S01--S05): conventional relay at $\fconv=20$\,Hz false-alarms
due to $\sigman(\fconv) \approx A_{\rm noise}$.  Optimal relay selects
$\fopt \approx 27$--$31$\,Hz where $\sigman \approx 0$.

Fault scenarios (S06--S10): high-resistance faults ($R_f \geq 6$\,pu) give
$\Ifault \leq 0.152$\,pu (margin $\leq 0.052$\,pu over $\Itrip$), erased by
$2\sigman(20\,\mathrm{Hz}) \geq 0.12$\,pu at conventional frequency.

\subsection{Monte Carlo parameters}

\begin{center}
\begin{tabular}{lll}
  \toprule
  Parameter & Distribution & Range \\
  \midrule
  Trials & — & 2\,000 (1\,000 security / 1\,000 fault) \\
  $\fsub$ & Uniform & [15, 25]\,Hz \\
  $A_{\rm noise}$ & Uniform & [0.03, 0.09]\,pu \\
  $R_f$ & Uniform & [0.0, 7.0]\,pu (fault trials) \\
  \bottomrule
\end{tabular}
\end{center}

%% ────────────────────────────────────────────────────────────────────────────
\section{Results}
\label{sec:results}
%% ────────────────────────────────────────────────────────────────────────────

\subsection{Deterministic results}

All ten scenarios pass at $\fopt$, with the optimal frequency ranging from
26.1 to 31.3\,Hz across scenarios:

\begin{center}
\begin{tabular}{llcccc}
  \toprule
  ID & Type & $\fopt$ (Hz) & $2\sigman(\fconv)$ & $2\sigman(\fopt)$
     & Opt.\ result \\
  \midrule
  S01 & SEC & 28.4 & 0.1600 & 0.0031 & \textcolor{darkgreen}{PASS} \\
  S02 & SEC & 30.4 & 0.1281 & 0.0031 & \textcolor{darkgreen}{PASS} \\
  S03 & SEC & 26.5 & 0.1281 & 0.0030 & \textcolor{darkgreen}{PASS} \\
  S04 & SEC & 31.3 & 0.0849 & 0.0031 & \textcolor{darkgreen}{PASS} \\
  S05 & SEC & 28.7 & 0.2000 & 0.0029 & \textcolor{darkgreen}{PASS} \\
  S06 & FLT & 28.4 & 0.1600 & 0.0031 & \textcolor{darkgreen}{PASS} \\
  S07 & FLT & 28.4 & 0.1600 & 0.0031 & \textcolor{darkgreen}{PASS} \\
  S08 & FLT & 30.4 & 0.1281 & 0.0031 & \textcolor{darkgreen}{PASS} \\
  S09 & FLT & 26.1 & 0.0961 & 0.0032 & \textcolor{darkgreen}{PASS} \\
  S10 & FLT & 28.6 & 0.1800 & 0.0030 & \textcolor{darkgreen}{PASS} \\
  \midrule
  \multicolumn{5}{l}{Total (optimal)} & 10/10 \\
  \multicolumn{5}{l}{Total (conventional $\fconv=20$\,Hz)} & 0/10 \\
  \bottomrule
\end{tabular}
\end{center}

\noindent The $50\times$ reduction in noise (from $2\sigman \approx 0.16$ to
$\approx 0.003$\,pu) is the key mechanism enabling reliable detection at
the same $\Itrip$ threshold.

\subsection{Monte Carlo results}
\label{sec:mc_results}

\begin{center}
\begin{tabular}{lcccc}
  \toprule
  Metric & Conventional & Optimal & Target & Status \\
  \midrule
  $P_D$ (dependability) & 0.813 & 1.000 & $\geq 0.970$
    & \textcolor{darkgreen}{\textbf{PASS}} \\
  $P_{\rm FA}$ (false alarm) & 0.819 & 0.000 & $\leq 0.010$
    & \textcolor{darkgreen}{\textbf{PASS}} \\
  $\fopt$ mean & — & 27.2\,Hz & $\in [25, 36]$\,Hz
    & \textcolor{darkgreen}{\textbf{PASS}} \\
  $\fopt$ std  & — & 3.3\,Hz  & —
    & — \\
  \bottomrule
\end{tabular}
\end{center}

\begin{remark}[False alarm improvement]
  The conventional relay at 20\,Hz has $P_{\rm FA} = 0.819$ under the sampled
  noise conditions ($\fsub \sim U[15,25]$\,Hz, $A_{\rm noise} \sim U[0.03,0.09]$
  pu).  The optimal relay achieves $P_{\rm FA} = 0.000$ --- an $81.9$ percentage
  point improvement --- because $\fopt$ is always shifted sufficiently far from
  $\fsub$ that $2\sigman(\fopt) \ll \Itrip - \Ihealthy(\fopt)$.
\end{remark}

\subsection{Frequency distribution}

The empirical distribution of $\fopt$ across MC trials (mean 27.2\,Hz, std
3.3\,Hz) confirms that the optimal frequency lies consistently above $\fsub$
and below the security crossover $f_{\rm up} \approx 35$\,Hz, in line with
Proposition~\ref{prop:fopt}.

%% ────────────────────────────────────────────────────────────────────────────
\section{Sensitivity Analysis}
\label{sec:sensitivity}
%% ────────────────────────────────────────────────────────────────────────────

\subsection{High-resistance fault detectability}

The maximum detectable fault resistance is
$R_{f,\rm max} = V_{\rm inj}/\Itrip - \RN = 9.92$\,pu (corresponding to
$\Ifault = \Itrip$).  For practical IBR networks with $A_{\rm noise} \leq 0.09$\,pu
and $\fopt \geq 25$\,Hz, the noise at $\fopt$ is $\sigman(\fopt) \leq 0.003$\,pu,
reducing the effective sensitivity limit to:
\[
  R_{f,\rm eff} = V_{\rm inj}/(\Itrip + 2\sigman(\fopt)) - \RN \approx 9.86\;\text{pu}
\]
--- a degradation of only 0.6\% compared to the noise-free case.

\subsection{IBR noise bandwidth}

The analysis assumes $\sigma_{\rm sub} = 3$\,Hz.  For wider noise bandwidths
($\sigma_{\rm sub}$ up to 5\,Hz), the required frequency shift to achieve
$\sigman(\fopt) \leq 0.003$\,pu increases from $\sim$8\,Hz to $\sim$14\,Hz,
keeping $\fopt$ within the valid range $[f_{\rm low}, f_{\rm up}]$ for
$\fsub \leq 22$\,Hz.

%% ────────────────────────────────────────────────────────────────────────────
\section{Comparison with Existing Methods}
\label{sec:comparison}
%% ────────────────────────────────────────────────────────────────────────────

\begin{center}
\begin{tabular}{p{0.30\textwidth}ccp{0.24\textwidth}}
  \toprule
  Method & $P_D$ & $P_{\rm FA}$ & Limitation \\
  \midrule
  Standard 64G ($f=20$\,Hz) \cite{IEEE_C37_101}
    & 0.813 & 0.819 & Fixed $f$ at noise peak \\
  DSP-filtered 64G \cite{Finney1988}
    & $\sim$0.90 & $\sim$0.10 & Filter delay; tuned to 20\,Hz only \\
  Proposed adaptive $\fopt$
    & 1.000 & 0.000 & Requires online $\fsub$ estimation \\
  \bottomrule
\end{tabular}
\end{center}

%% ────────────────────────────────────────────────────────────────────────────
\section{Implementation Notes}
\label{sec:impl}
%% ────────────────────────────────────────────────────────────────────────────

The adaptive injection frequency scheme requires:
\begin{enumerate}
  \item \textbf{IBR noise peak estimation}: a sliding-window FFT or MUSIC
        spectral estimator to track $\hat{f}_{\rm sub}$ in real time.
        Latency $\leq 100$\,ms is achievable with a 0.5\,s window.
  \item \textbf{Frequency-agile injection source}: the neutral injection
        transformer must support $f \in [20, 40]$\,Hz.  Modern DSP-controlled
        injection units accommodate this range with no hardware change.
  \item \textbf{Trip threshold recalculation}: $\Itrip$ is fixed; only the
        injection frequency changes.  The relay logic is otherwise identical to
        standard 64G.
\end{enumerate}

%% ────────────────────────────────────────────────────────────────────────────
\section{Conclusions}
\label{sec:conclusions}
%% ────────────────────────────────────────────────────────────────────────────

\begin{enumerate}
  \item Fixed-frequency (20\,Hz) Relay 64G protection fails catastrophically
        under IBR sub-harmonic noise: $P_{\rm FA} = 0.819$ and $P_D = 0.813$
        at conventional settings.

  \item An adaptive injection frequency $\fopt$ that maximises the balanced
        security/dependability margin eliminates both failure modes:
        $P_D = 1.000$, $P_{\rm FA} = 0.000$.

  \item The fault current $\Ifault(R_f)$ is frequency-independent (conductive
        path), so optimal frequency selection reduces only the noise term
        $2\sigman(\fopt)$, not the signal.  This decoupling enables a clean
        analytical bound (Proposition~\ref{prop:fopt}).

  \item $\fopt$ lies consistently in $[25, 36]$\,Hz for $\fsub \in [15, 25]$\,Hz,
        providing a design rule: \emph{inject at least 8\,Hz above the
        measured IBR noise peak}.

  \item The scheme is backward-compatible with existing 64G relay hardware;
        only the injection source frequency control requires modification.
\end{enumerate}

%% ────────────────────────────────────────────────────────────────────────────
\appendix
\section{Simulation Files}
\label{app:files}
%% ────────────────────────────────────────────────────────────────────────────

\begin{tabular}{p{0.52\textwidth}p{0.42\textwidth}}
  \toprule
  File & Description \\
  \midrule
  \texttt{relay64g/run\_tr61\_relay64g\_subharm.py}
    & TR-61 simulation driver: physics model, optimal-$f$ solver, 10-scenario + MC \\
  \texttt{relay64g/outputs/tr61/tr61\_deterministic.csv}
    & 10-scenario results (10/10 PASS) \\
  \texttt{relay64g/outputs/tr61/tr61\_mc\_metrics.csv}
    & Monte Carlo summary ($P_D$, $P_{\rm FA}$, $\fopt$ statistics) \\
  \texttt{TR61\_Relay64G\_SubHarmonic/main\_report61.tex}
    & This document \\
  \bottomrule
\end{tabular}

\bibliographystyle{IEEEtran}
\bibliography{references61}

\end{document}
